\centerline{\bf \Huge Сравнения по модулю многочленов}

\begin{flushright}
        \scriptsize{Всё таки первый враг человека --- это он сам\\
                    Гендо Икари}
\end{flushright}

\problem{1}
При каких натуральных $n$ многочлен $(x+1)^n+x^n+1$ делится на $x^2+x+1$ ?

\problem{2}
В лес за грибами пошли 11 девочек и $n$ мальчиков. Вместе они собрали $n^2+9 n-2$ гриба, причём все они собрали поровну грибов. Кого было больше: мальчиков или девочек?

\problem{3}
$P(x)$ и $Q(x)$ - многочлены с действительными коэффициентами. Известно, что многочлен $P\left(x^3\right)+x Q\left(x^3\right)$ делится на $x^3-1$. Докажите, что $P(x)$ и $Q(x)$ оба делятся на $x-1$.

\problem{4}
Докажите, что для любого многочлена $P(x)$ многочлен $P(P(\ldots P(x) \ldots))-x$ делится на $P(x)-x$.

\problem{5}
Докажите, что для любого многочлена $P$ с целыми коэффициентами и любого натурального $k$ существует такое натуральное $n$, что $P(1)+P(2)+\ldots+P(n)$ делится на $k$.

\problem{6}
При каких натуральных $k$ многочлен $1+x^4+x^8+\ldots+x^{4 k}$ делится на многочлен $1+x^2+x^4+\ldots+x^{2 k}$ ?

\problem{7}
Найдите все пары натуральных чисел $x, y$ такие, что $x^3+1$ делит $(x+1)^y$.

\problem{8}
Найдите все многочлены $P(x)$ с целыми коэффициентами такие, что $P(P(x)+x)$ является простым числом при бесконечном количестве натуральных $x$.